\chapter{Stroke}

\section*{Patient focus}
When someone has stroke symptoms, a CT scan is done first and very quickly. The main reason is to determine whether the stroke is caused by a blocked blood vessel (ischaemic) or a burst blood vessel (haemorrhagic). This distinction is critical because the treatments are completely different. A blocked vessel may be treated with clot-busting drugs or surgery, while a bleed requires different management. Time is brain---the faster imaging is done, the more brain can be saved.


\label{ch:ch60_stroke}

\section{Acute Assessment}
Stroke imaging is time-critical. The target is to complete non-contrast CT within \textbf{20 minutes} of arrival for patients presenting within the thrombolysis window.

\subsection{Initial Imaging}
\begin{itemize}
  \item \textbf{Non-contrast CT (NCCT)}: first-line to exclude haemorrhage. Early ischaemic changes may be subtle---sulcal effacement, loss of grey-white differentiation, hyperdense vessel sign (acute thrombus)
  \item \textbf{CT angiography (CTA)}: from aortic arch to vertex to identify large vessel occlusion (LVO). Targets: internal carotid, M1/M2 MCA, basilar
  \item \textbf{CT perfusion (CTP)}: identifies the ischaemic core (irreversibly damaged) and penumbra (potentially salvageable tissue)
\end{itemize}

\subsection{MRI for Stroke}
MRI is more sensitive for early ischaemia and posterior circulation strokes:
\begin{itemize}
  \item \textbf{DWI}: often shows restricted diffusion early and is highly sensitive for acute infarction, although a negative study does not exclude every acute stroke.
  \item \textbf{FLAIR}: helps characterize lesion age and distribution; DWI-FLAIR mismatch can support selection for imaging-based treatment in appropriate protocols, but is not a precise clock for stroke onset.
  \item \textbf{MRA}: non-invasive vessel assessment
  \item \textbf{SWI/T2*}: detects microbleeds and haemorrhagic transformation
\end{itemize}

\section{Ischaemic vs Haemorrhagic Stroke}

\begin{table}[h]
\centering
\caption{Imaging features distinguishing stroke types}
\begin{tabular}{|l|l|l|}
\hline
\textbf{Feature} & \textbf{Ischaemic} & \textbf{Haemorrhagic} \\
\hline
CTA & Vessel occlusion may be seen & Usually no occlusion \\
DWI & Bright (restricted diffusion) & Variable \\
Early CT signs & Sulcal effacement, loss of G/W differentiation & Immediate hyperdensity \\
NCCT appearance & Hypodense area (late) & Hyperdense lesion \\
Vascular territory & Corresponds to artery & Usually deep structures or lobar \\
\hline
\end{tabular}
\end{table}

\section{ASPECTS Score}
The \textbf{Alberta Stroke Programme Early CT Score} quantifies early ischaemic change on NCCT in MCA territory strokes:

\begin{itemize}
  \item Score 10 = normal brain
  \item 10 regions assessed (caudate, lentiform, internal capsule, insular ribbon, MCA cortical territories M1--M6)
  \item Each region with early ischaemic change loses 1 point
  \item Score $\leq$7: large core infarct, less benefit from recanalisation
  \item Score 8--10: small core, good candidate for intervention
\end{itemize}

\begin{pearlbox}[ASPECTS and treatment decisions]
ASPECTS $<$6 generally indicates extensive established infarction where thrombectomy benefit is reduced but not eliminated. ASPECTS 8--10 with LVO is the ideal candidate for thrombectomy. Report ASPECTS promptly---it directly influences the interventional decision.
\end{pearlbox}

\section{The Mismatch Concept}
The core--penumbra mismatch is the basis for extended-window thrombolysis and thrombectomy decisions:

\begin{itemize}
  \item \textbf{Core}: irreversibly infarcted tissue (CBF $<$30\% of normal on CTP, or bright on DWI)
  \item \textbf{Penumbra}: hypoperfused but viable tissue (Tmax $>$6 s on CTP)
  \item \textbf{Mismatch}: penumbra $>$ core by $\geq$20\% of core volume, and absolute mismatch $>$1.5 mL
  \item \textbf{Target mismatch}: favourable profile suggesting salvageable tissue
\end{itemize}

\textbf{Clinical application}:
\begin{enumerate}
  \item \textbf{$<$4.5 hours}: NCCT to exclude haemorrhage, then thrombolysis (IV tPA)---CTP not required
  \item \textbf{4.5--24 hours}: CTP or MRI with DWI/PWI to identify mismatch; thrombectomy for LVO with target mismatch
  \item \textbf{Wake-up stroke}: DWI-FLAIR mismatch on MRI (DWI positive, FLAIR negative) suggests stroke within 4.5 hours; treat with thrombolysis
\end{enumerate}

\section{Vessel Imaging}
CTA or MRA identifies:
\begin{itemize}
  \item \textbf{Large vessel occlusion (LVO)}: ICA, M1, basilar---amenable to thrombectomy
  \item \textbf{Intracranial stenosis}: May influence treatment and secondary prevention
  \item \textbf{Dissection}: intimal flap, string sign, aneurysmal dilation---common in young stroke
  \item \textbf{Collaterals}: good collaterals predict better outcomes and slower infarct growth
\end{itemize}

\section{When to Recommend Thrombolysis vs Thrombectomy}

\begin{table}[h]
\centering
\caption{Treatment decision based on imaging}
\begin{tabular}{|l|l|l|}
\hline
\textbf{Scenario} & \textbf{Treatment} & \textbf{Imaging requirement} \\
\hline
4.5--24 h, LVO, mismatch & Thrombectomy & CTA + CTP or MRI \\
$<$4.5 h, no haemorrhage & IV thrombolysis & NCCT only \\
Basilar occlusion & Thrombectomy & CTA \\
Large core ($>$70 mL), no mismatch & Medical management & CTP/MRI \\
Wake-up stroke, DWI/FLAIR mismatch & IV thrombolysis & MRI \\
\hline
\end{tabular}
\end{table}

\section{Posterior Circulation Stroke}
Posterior circulation strokes are more challenging to diagnose:
\begin{itemize}
  \item NCCT is often normal early---use MRI with DWI
  \item Basilar artery occlusion on CTA is an emergency---thrombectomy even beyond 24 hours in selected patients
  \item Posterior fossa strokes cause wall-eyed bilateral syndrome (WEBINO), lateral medullary syndrome, and cerebellar infarction with mass effect
  \item CT perfusion may be less reliable in the posterior fossa due to beam-hardening artefact
\end{itemize}

\begin{pearlbox}[Don't miss posterior circulation stroke]
Persistent vertigo, diplopia, dysarthria, ataxia, or crossed signs demand MRI with DWI even if NCCT is normal. A normal CT does not exclude posterior circulation infarction.
\end{pearlbox}

\section{Complications of Stroke on Imaging}
\begin{itemize}
  \item \textbf{Haemorrhagic transformation}: petechial (Type 1--2) or space-occupying (Type 3--4)---Type 4 requires stopping anticoagulation
  \item \textbf{Malignant MCA infarction}: large territorial infarction with midline shift---may need decompressive craniectomy within 48 hours
  \item \textbf{Cerebral oedema}: peaks at 3--5 days; mass effect, sulcal effacement, ventricular compression
  \item \textbf{Vasogenic oedema around tumour}: if stroke-like presentation is actually tumour---look for enhancement and mass effect
\end{itemize}

\section{Safety}
Imaging must be coordinated with last-known-well time, neurologic examination, glucose, anticoagulation, renal function, pregnancy, and treatment eligibility. Do not delay time-sensitive therapy for nonessential sequences.

\clearpage
